\section{Projective modules}
\label{section-projective}

\noindent
Some lemmas on projective modules.

\begin{definition}
\label{definition-projective}
Let $R$ be a ring. An $R$-module $P$ is {\it projective} if and only if
the functor $\Hom_R(P, -) : \text{Mod}_R \to \text{Mod}_R$ is
an exact functor.
\end{definition}

\noindent
The functor $\Hom_R(M, - )$ is left exact for any $R$-module $M$, see
Lemma \ref{lemma-hom-exact}.
Hence the condition for $P$ to be projective really signifies that given
a surjection of $R$-modules $N \to N'$ the map
$\Hom_R(P, N) \to \Hom_R(P, N')$ is surjective.

\begin{lemma}
\label{lemma-characterize-projective}
Let $R$ be a ring. Let $P$ be an $R$-module.
The following are equivalent
\begin{enumerate}
\item $P$ is projective,
\item $P$ is a direct summand of a free $R$-module, and
\item $\Ext^1_R(P, M) = 0$ for every $R$-module $M$.
\end{enumerate}
\end{lemma}

\begin{proof}
Assume $P$ is projective. Choose a surjection $\pi : F \to P$ where $F$
is a free $R$-module. As $P$ is projective there exists a
$i \in \Hom_R(P, F)$ such that $\pi \circ i = \text{id}_P$.
In other words $F \cong \Ker(\pi) \oplus i(P)$ and we see
that $P$ is a direct summand of $F$.

\medskip\noindent
Conversely, assume that $P \oplus Q = F$ is a free $R$-module.
Note that the free module $F = \bigoplus_{i \in I} R$ is projective
as $\Hom_R(F, M) = \prod_{i \in I} M$ and the functor
$M \mapsto \prod_{i \in I} M$ is exact.
Then $\Hom_R(F, -) = \Hom_R(P, -) \times \Hom_R(Q, -)$
as functors, hence both $P$ and $Q$ are projective.

\medskip\noindent
Assume $P \oplus Q = F$ is a free $R$-module. Then we have a
free resolution $F_\bullet$ of the form
$$
\ldots F \xrightarrow{a} F \xrightarrow{b} F \to P \to 0
$$
where the maps $a, b$ alternate and are equal to the projector onto
$P$ and $Q$. Hence the complex $\Hom_R(F_\bullet, M)$ is split
exact in degrees $\geq 1$, whence we see the vanishing in (3).

\medskip\noindent
Assume $\Ext^1_R(P, M) = 0$ for every $R$-module $M$.
Pick a free resolution $F_\bullet \to P$. Set
$M = \Im(F_1 \to F_0) = \Ker(F_0 \to P)$.
Consider the element $\xi \in \Ext^1_R(P, M)$ given by
the class of the quotient map $\pi : F_1 \to M$. Since $\xi$ is zero
there exists a map $s : F_0 \to M$ such that $\pi = s \circ (F_1 \to F_0)$.
Clearly, this means that
$$
F_0 = \Ker(s) \oplus \Ker(F_0 \to P) =
P \oplus \Ker(F_0 \to P)
$$
and we win.
\end{proof}

\begin{lemma}
\label{lemma-characterize-finite-projective-noetherian}
Let $R$ be a Noetherian ring. Let $P$ be a finite $R$-module.
If $\Ext^1_R(P, M) = 0$ for every finite $R$-module $M$, then
$P$ is projective.
\end{lemma}

\noindent
This lemma can be strengthened: There is a version for finitely presented
$R$-modules if $R$ is not assumed Noetherian. There is a version with $M$
running through all finite length modules in the Noetherian case.

\begin{proof}
Choose a surjection $R^{\oplus n} \to P$ with kernel $M$.
Since $\Ext^1_R(P, M) = 0$ this surjection is split and
we conclude by Lemma \ref{lemma-characterize-projective}.
\end{proof}

\begin{lemma}
\label{lemma-direct-sum-projective}
A direct sum of projective modules is projective.
\end{lemma}

\begin{proof}
This is true by the characterization of projectives as direct
summands of free modules in
Lemma \ref{lemma-characterize-projective}.
\end{proof}

\begin{lemma}
\label{lemma-lift-projective-module}
Let $R$ be a ring. Let $I \subset R$ be a nilpotent ideal. Let
$\overline{P}$ be a projective $R/I$-module. Then there exists a
projective $R$-module $P$ such that $P/IP \cong \overline{P}$.
\end{lemma}

\begin{proof}
By Lemma \ref{lemma-characterize-projective}
we can choose a set $A$ and a direct sum decomposition
$\bigoplus_{\alpha \in A} R/I = \overline{P} \oplus \overline{K}$
for some $R/I$-module $\overline{K}$. Write $F = \bigoplus_{\alpha \in A} R$
for the free $R$-module on $A$. Choose a lift
$p : F \to F$ of the projector $\overline{p}$ associated to
the direct summand $\overline{P}$ of $\bigoplus_{\alpha \in A} R/I$.
Note that $p^2 - p \in \text{End}_R(F)$ is a nilpotent
endomorphism of $F$ (as $I$ is nilpotent and the matrix entries of
$p^2 - p$ are in $I$; more precisely, if $I^n = 0$, then $(p^2 - p)^n = 0$).
Hence by Lemma \ref{lemma-lift-idempotents-noncommutative}
we can modify our choice of $p$ and assume that $p$ is a projector.
Set $P = \Im(p)$.
\end{proof}

\begin{lemma}
\label{lemma-lift-finite-projective-module}
Let $R$ be a ring. Let $I \subset R$ be a locally nilpotent ideal. Let
$\overline{P}$ be a finite projective $R/I$-module. Then there exists a
finite projective $R$-module $P$ such that $P/IP \cong \overline{P}$.
\end{lemma}

\begin{proof}
Recall that $\overline{P}$ is a direct summand of a free $R/I$-module
$\bigoplus_{\alpha \in A} R/I$ by Lemma \ref{lemma-characterize-projective}.
As $\overline{P}$ is finite, it follows that $\overline{P}$ is contained
in $\bigoplus_{\alpha \in A'} R/I$ for some $A' \subset A$ finite.
Hence we may assume we have a direct sum decomposition
$(R/I)^{\oplus n} = \overline{P} \oplus \overline{K}$
for some $n$ and some $R/I$-module $\overline{K}$. Choose a lift
$p \in \text{Mat}(n \times n, R)$ of the projector $\overline{p}$
associated to the direct summand $\overline{P}$ of $(R/I)^{\oplus n}$.
Note that $p^2 - p \in \text{Mat}(n \times n, R)$ is nilpotent:
as $I$ is locally nilpotent and the matrix entries $c_{ij}$ of
$p^2 - p$ are in $I$ we have $c_{ij}^t = 0$ for some $t > 0$ and
then $(p^2 - p)^{tn^2} = 0$ (by looking at the matrix coefficients).
Hence by Lemma \ref{lemma-lift-idempotents-noncommutative}
we can modify our choice of $p$ and assume that $p$ is a projector.
Set $P = \Im(p)$.
\end{proof}

\begin{lemma}
\label{lemma-lift-projective}
Let $R$ be a ring.
Let $I \subset R$ be an ideal.
Let $M$ be an $R$-module.
Assume
\begin{enumerate}
\item $I$ is nilpotent,
\item $M/IM$ is a projective $R/I$-module,
\item $M$ is a flat $R$-module.
\end{enumerate}
Then $M$ is a projective $R$-module.
\end{lemma}

\begin{proof}
By Lemma \ref{lemma-lift-projective-module} we can find a projective
$R$-module $P$ and an isomorphism $P/IP \to M/IM$. We are going to show
that $M$ is isomorphic to $P$ which will finish the proof. Because $P$
is projective we can lift the map $P \to P/IP \to M/IM$ to an $R$-module
map $P \to M$ which is an isomorphism modulo $I$. Since $I^n = 0$
for some $n$, we can use the filtrations
\begin{align*}
0 = I^nM \subset I^{n - 1}M \subset \ldots \subset IM \subset M \\
0 = I^nP \subset I^{n - 1}P \subset \ldots \subset IP \subset P
\end{align*}
to see that it suffices to show that the induced maps
$I^aP/I^{a + 1}P \to I^aM/I^{a + 1}M$ are bijective. Since both $P$
and $M$ are flat $R$-modules we can identify this with the map
$$
I^a/I^{a + 1} \otimes_{R/I} P/IP
\longrightarrow
I^a/I^{a + 1} \otimes_{R/I} M/IM
$$
induced by $P \to M$. Since we chose $P \to M$ such that the induced
map $P/IP \to M/IM$ is an isomorphism, we win.
\end{proof}

\begin{lemma}
\label{lemma-projective-modulo-two-ideals}
Let $R$ be a ring. Let $I, J \subset R$ be ideals such that $I \cap J = 0$.
Let $P$ be an $R$-module such that $P/IP$ is a projective $R/I$-module and
$P/JP$ is a projective $R/J$-module. Then $P$ is a projective $R$-module.
\end{lemma}

\begin{proof}
Choose a surjection $p : F \to P$ where $F$ is a free $R$-module.
Since $P/IP$ is a projective $R/I$-module, we can choose a map
$f : P/IP \to F/IF$ which is a right inverse to $p \bmod I$.
Consider the map $q : F/JF \to F/(I + J)F \times_{P/(I + J)P} P/JP$
induced by the quotient map $F/JF \to F/(I + J)F$ and $p$.
Note that $q$ is a surjective map of $R/J$-modules (small detail omitted).
Consider the map $f' : P/JP \to F/(I + J)F \times_{P/(I + J)P} P/JP$
induced by $f$ and the identity map.
Since $P/JP$ is a projective $R/J$-module and $q$ surjective,
we can choose a map $g : P/JP \to F/JF$ such that $q \circ g = f'$.
Then $f \bmod I + J = g \bmod I + J$. Since $F = F/JF \times_{F/(I + J)F} F/IF$
because $I \cap J = 0$, we conclude that we obtain a well defined homomorphism
$h = (f, g) : P \to F$ of $R$-modules. The map $a = p \circ h : P \to P$
reduces to the identity modulo $I$ and modulo $J$. Then $a$ also induces
the identity map on the submodule $IP$: if $x = \sum i_\alpha x_\alpha$
with $i_\alpha \in I$ and $x_\alpha \in P$, then
$a(x) = \sum i_\alpha a(x_\alpha) = \sum i_\alpha x_\alpha = x$
because $a(x_\alpha) = x_\alpha + j_\alpha$ with $j_\alpha \in J$
and $i_\alpha j_\alpha = 0$. Thus $a : P \to P$ acts as the identity
on $IP$ and on $P/IP$ and we conclude that $a$ is an automorphism of $P$
(small detail omitted). Thus $P$ is a summand of $F$, whence projective.
\end{proof}










